%%%%%%%%%%%%%%%%%%%%%%%%%%%%%%%%%%%%%%%%%%%%%%%%%%%%%%%%%%%%%%%%%%%%%%%%%%%%%
%% INÍCIO CAPÍTULO                                                         %%
%%%%%%%%%%%%%%%%%%%%%%%%%%%%%%%%%%%%%%%%%%%%%%%%%%%%%%%%%%%%%%%%%%%%%%%%%%%%%

\chapter{Conclusões}
\label{cap:4:conclusoes}

Este trabalho apresentou o desenvolvimento de um \textit{boilerplate} para a criação de aplicações e ferramentas interativas destinadas à análise e visualização de dados. A solução proposta combina ferramentas amplamente utilizadas, como \texttt{GTK}, \texttt{matplotlib}, \texttt{pandas} e \texttt{scikit-learn}, em uma arquitetura modular e extensível. Neste capítulo, é apresentado os resultados obtidos e listado possíveis trabalhos futuros que podem expandir ou aprimorar o \textit{boilerplate} inicial.

\section{Resultados}

O principal resultado deste trabalho foi o desenvolvimento de uma estrutura base que simplifica a criação de ferramentas de análise e visualização de dados. Entre as principais funcionalidades, destacam-se:

\begin{itemize}
    \item Interface modular e extensível: A utilização do GTK para a criação da interface gráfica permitiu a construção de dashboards interativos, organizados em módulos independentes.
    
    \item Visualização eficiente de dados: A integração com o \texttt{matplotlib} e \texttt{seaborn} garantiu uma visualização clara e customizável, possibilitando análises mais intuitivas para os usuários.
    
    \item Gerenciamento genérico de modelos: O gerenciador de modelo baseado no \texttt{FAModel} demonstrou ser eficaz para lidar com algoritmos de processamento de dados e aprendizado de máquina, como o \texttt{KMeans}, permitindo armazenamento, carregamento e reutilização de modelos de maneira consistente e eficiente.
    
    \item Escalabilidade e reutilização: A arquitetura orientada a objetos proporcionou modularidade e reutilização de componentes, facilitando a extensão do \textit{boilerplate} para novos algoritmos e funcionalidades.

    \item Documentação: A integração com \texttt{mkdocs} permiti a criação de uma documentação clara e acessível, facilitando a documentação de experimentos, resultados e funcionalidades do sistema quando estendido.

    \item \textit{linter} e \textit{formatter}: A integração com \texttt{ruff} entrega uma ferramenta de formatação e verificação de código, garantindo a consistência e a qualidade do código produzido.

    \item Compartilhamento de arquivos JSON e modelos binários: A capacidade de exportar e importar dados e modelos em formatos comuns, como JSON e \texttt{pickle}, facilita a colaboração e a reutilização de resultados.
\end{itemize}

Os exemplos fornecidos demonstraram que o \textit{boilerplate} é capaz de integrar diferentes algoritmos de processamento de dados e aprendizado de máquina, de forma simples e oferecer suporte a fluxos completos de análise de dados, desde a manipulação de entradas até a geração de resultados visuais.

\section{Trabalhos Futuros}

Apesar dos avanços alcançados, este trabalho abre espaço para diversas possibilidades de melhorias e expansões:

\begin{itemize}
    \item Integração com outros frameworks: Explorar a integração com outras bibliotecas de visualização e aprendizado de máquina, como \texttt{Plotly} para gráficos interativos ou gerenciadores específicos para \texttt{TensorFlow} para modelos mais complexos.

    \item Processamento distribuído: Implementar suporte para processamento distribuído, permitindo a execução de análises em paralelo em múltiplos núcleos ou máquinas.
    
    \item Otimização de desempenho: Implementar melhorias para reduzir o consumo de memória ao lidar com grandes volumes de dados, como a utilização de \textit{generators} e \textit{streams}.

    \item Documentação expandida: Ampliar a documentação para incluir tutoriais, exemplos e guias de uso, facilitando a adoção e a extensão do sistema a partir do próprio repositório através do \texttt{mkdocs}. Essa estratégia também fornecerá exemplos de como documentar as próprias aplicações desenvolvidas com o \textit{boilerplate}.

    \item Traduzir e internacionalizar: Adicionar suporte a outros idiomas e localizações, como espanhol, japonês e português, tornando o sistema acessível a um público mais amplo.

    \item Processamento em tempo real: Adicionar suporte para processamento em tempo real, permitindo a análise de dados em fluxo contínuo, como em aplicações de IoT ou monitoramento de sistemas.

    \item Suporte a múltiplas fontes de dados facilitado: Embora a biblioteca \texttt{pandas} já forneça conectores, a proposta de melhoria é expandir o sistema para suportar a integração com diferentes fontes de dados, como bancos de dados SQL, NoSQL ou APIs de serviços web, fornecendo os widgets necessários para interação com essas fontes.
    
    \item Interface web: Expandir a aplicação para suportar sistemas de visualização baseados na web, ampliando o alcance e a usabilidade do sistema.
    
    \item Automação de análises: Incorporar ferramentas para automação de análises preditivas e geração de relatórios, tornando o sistema ainda mais útil em contextos empresariais.
    
    \item Validação com usuários: Realizar testes de usabilidade com potenciais usuários para identificar melhorias na experiência e na interface do sistema.
    
    \item Biblioteca de modelos: Desenvolver uma biblioteca ampliada de algoritmos prontos para uso, abrangendo métodos de regressão, classificação e análise de séries temporais.
\end{itemize}

Essas direções apontam para a evolução do \textit{boilerplate} como uma solução completa e flexível para análise de dados, tanto em contextos acadêmicos quanto industriais.

%%%%%%%%%%%%%%%%%%%%%%%%%%%%%%%%%%%%%%%%%%%%%%%%%%%%%%%%%%%%%%%%%%%%%%%%%%%%%
%% FIM CAPÍTULO                                                            %%
%%%%%%%%%%%%%%%%%%%%%%%%%%%%%%%%%%%%%%%%%%%%%%%%%%%%%%%%%%%%%%%%%%%%%%%%%%%%%
